% Gráficos canónicos de D3 en su orientación original.
\ProvideDocumentCommand\figuraHorizontalAqui{mmm}{\par\figuraHorizontal{#1}{#2}{#3}\clearpage}
% Contenido de D3 trasladado desde sus Markdown del Informe 2.
La gestión combina PMBOK para alcance, adquisiciones, riesgos y cambios; Scrum para incrementos de software; Kanban para incidencias y estabilización. S3/T-12 fija 28 RF y 14 RNF —catálogo de la oferta, no número de requisitos impuestos por el Caso—; S7 organiza paquetes e hitos y S9 establece aceptación y evidencia. Los Formularios T-9/T-10 acompañan el capítulo como archivos propios conforme al Comunicado 10.

El CLIENTE recibe una línea base controlada, responsables y decisiones registradas, incrementos trazados a requisitos y un flujo DevSecOps con bloqueo de promoción ante incumplimiento. El contrato conserva 56 meses, con desarrollo E1 M1–12, marcha blanca E1 M13–15, producción E1 M16, desarrollo E2 M13–18, marcha blanca E2 M19–20 y operación completa M21–56 (\art{17.1}{12}). La metodología no acredita por sí sola que esos resultados ya se hayan obtenido.

La infraestructura y el montaje requieren planificación predictiva por terminal y unidad; las interfaces, contratos y modelos de costo requieren validación iterativa con usuarios; los incidentes 24x7 requieren flujo con prioridad y escalamiento. Los frentes se sincronizan mediante los comités contractuales y las puertas de S9, sin alterar hitos obligatorios.

\section{Metodología de Gestión de Proyectos}\label{sec:6-gestion}\label{sec:06-metodologia-de-gestion-del-proyecto}

La gestión del proyecto adopta las directrices de la guía PMBOK \parencite{pmbok7}, adaptada específicamente a las restricciones operativas y contractuales del transporte de carga por carretera y la logística interurbana (\art{4.3}{5}, \transv{RT-19.01}{33}). Este marco mantiene trazabilidad verificable sobre la línea base de alcance, tiempo y costo, articulando la relación contractual con Transportes Curimón S.A. mediante canales formales de rendición de cuentas.

\FloatBarrier
\subsection{Marco híbrido de gestión: PMBOK y enfoques ágiles}\label{sec:6-gestion-pmbok}\label{sec:06-marco-hibrido-de-gestion-pmbok-y-enfoques-agiles}

La administración del contrato se fundamenta en un modelo de ciclo de vida híbrido, el cual combina procesos predictivos para los compromisos contractuales vinculantes con ciclos adaptativos para la construcción de los módulos de software. La justificación de este enfoque dual se resume en los siguientes ejes operativos:

\begin{itemize}
\item Capa Predictiva (PMBOK 7): Aplica a la gestión de la Ruta Crítica contractual del Artículo 17°, la entrega de hitos del Formulario E-25, la integración del parque de \textbf{374} camiones, distinguiendo equipos existentes y las \textbf{182} instalaciones previstas en T-11 (\caso{numeral 14.1}{29}), las adecuaciones de salas técnicas y la obtención de recepciones provisorias y definitivas.
\item Capa Adaptativa (Scrum): Aplica al diseño y programación de los módulos aplicativos (torre de control, despacho, portal de transportistas, aplicación de conductores y analítica). Se estructura en iteraciones cortas (sprints de 2 semanas) con entregables potencialmente desplegables, lo que permite retroalimentación temprana de las contrapartes operativas de Curimón sin comprometer los plazos mayores de entrega.
\item Capa de Flujo Continuo (Kanban): Aplica a la atención de incidencias durante las marchas blancas contractuales: Etapa 1 en los meses 13–15 y Etapa 2 en los meses 19–20 (\art{17.1}{12}). También cubre la gestión de parches y solicitudes de servicio durante la operación de los meses 21–56.
\end{itemize}

\decision{titulo={Adopción de marco de gestión híbrido predictivo-ágil}, decide={Implementar gobernanza predictiva PMBOK para control contractual e hitos combinada con Scrum en iteraciones de dos semanas para software}, descarta={Enfoque puramente predictivo tipo cascada y enfoque puramente ágil sin línea base}, criterio={El transporte crítico exige certidumbre contractual de plazos de implantación junto con flexibilidad en diseño de interfaces y analítica}, fuente={\art{4.3}{5}, \transv{RT-19.01}{33}}}

La articulación temporal del marco de gestión y su integración con los hitos contractuales de la licitación se ilustra en la \figref{06-6-1-hitos}.

\figuraHorizontalAqui{figuras/06-metodologias/6-1-hitos.pdf}{Gobernanza temporal del proyecto y articulación de fases, etapas y marchas blancas}{06-6-1-hitos}

La \figref{06-6-1-hitos} representa el solapamiento desde M13: la estabilización de E1 y la construcción de E2 usan equipos y capacidad segregados. M25–56 aparece como tramo comprimido, identificado expresamente; el detalle mensual y las fechas derivadas de SA-01 corresponden a S7.

\FloatBarrier
\subsection{Gobernanza de interesados y plan de comunicaciones}\label{sec:6-gestion-interesados}\label{sec:06-gobernanza-de-interesados-y-plan-de-comunicaciones}

La gestión eficaz de los grupos de interés resulta determinante debido a la dispersión geográfica de las faenas y la multiplicidad de actores involucrados en la cadena logística. Se identifican y clasifican los siguientes actores clave:

\begin{itemize}
\item Patrocinador Ejecutivo y Dirección de Transporte Curimón: Enfocados en retorno de inversión, continuidad operacional del negocio y cumplimiento de contratos de flete minero, vitivinícola y retail.
\item Supervisores de Tráfico y Operadores de Torre de Control: Usuarios intensivos de la plataforma, orientados a la visibilidad en tiempo real, alertas de desvío y cumplimiento de itinerarios.
\item Conductores de Flota Propia y Terceros: \textbf{454} (\caso{numeral 14.1}{29}) conductores que interactúan directamente con la aplicación móvil y los sensores de cabina, priorizando la ergonomía, simplicidad y registro certero de jornadas laborales conforme al Artículo 25 bis del Código del Trabajo.
\item Transportistas Terceros y Dueños de Camiones: Propietarios de los 226 camiones subcontratados, interesados en la liquidación expedita de servicios y visibilidad telemática homologada.
\item Organismos Fiscalizadores (Dirección del Trabajo, MTT, SEC): Entidades que auditan la legalidad del transporte, pesos por eje, transporte de sustancias peligrosas y registros de jornada.
\end{itemize}

La lista anterior agrupa usuarios de operación; el plan de comunicaciones cubre, además, los trece grupos diferenciados en S2. La \tabref{06-actores-comunicacion} indica emisor, destinatario, información, cadencia y canal. Una solicitud urgente sigue el procedimiento de incidente o de cambio y no espera el informe periódico.

\begin{tabla}[ancho=completo, fuente=condensada]{Comunicaciones para los trece grupos de interés del diagnóstico}{06-actores-comunicacion}{L{27mm} X X X}
Destinatario & Información y frecuencia & Emisor y canal & Escalamiento \\ \midrule
Fondo de inversión & Riesgo financiero y continuidad; trimestral & Gerencia de Curimón; informe de directorio & Gerencia General \\
Dirección del Trabajo & Evidencia de jornada ante requerimiento formal & Encargado legal de Curimón; canal oficial & Gerencia y asesoría legal \\
Aseguradora de carga & Siniestros y cadena de evidencia; por evento & Riesgos de Curimón; expediente formal & Gerencia y Comité Ejecutivo \\
Familia fundadora & Hitos, continuidad y riesgos; mensual & Gerencia General; informe ejecutivo & Comité Ejecutivo \\
Gerencia General & Hitos, costos y riesgos; mensual & Jefe de Proyecto; Comité Ejecutivo & Administrador del Contrato \\
Operaciones & Despacho, bloqueos y cambios; semanal & Líder funcional; reunión y tablero & Comité de Proyecto \\
Finanzas & Conciliación, costo y liquidación; quincenal & Líder de datos; tablero y acta & Comité de Proyecto \\
Mantenimiento & Ventanas, kits y vigencias; semanal & Líder de implantación; agenda y acta & Comité de Operación \\
Prevención de Riesgos & Jornada, habilitaciones y excepciones; semanal & Líder funcional; registro de casos & Comité de Proyecto \\
TI y Control de Flota & Interfaces, seguridad y telemetría; semanal & Arquitecto; ADR y mesa técnica & Comité de Arquitectura \\
Conductores propios & Cambios de uso, privacidad y alertas; por liberación y turno & Capacitación; sesión y material por perfil & Supervisor de turno \\
Transportistas terceros & Adhesión, permisos y liquidación; por hito y mensualmente & Relación con terceros; portal y comunicación formal & Administrador del Contrato \\
Cliente exportador & Seguimiento y emisiones autorizadas; por corte contractual & Gerencia comercial; reporte aprobado & Gerencia General \\
\end{tabla}

La \tabref{06-actores-comunicacion} no concede al proponente comunicación directa en nombre del mandante con fiscalizadores, aseguradora o accionistas. Curimón emite esas comunicaciones; audIT prepara la evidencia que le corresponde. Todo acuerdo que cambie alcance, costo, plazo o nivel de servicio pasa por el Artículo 72 y queda en acta.

Conforme a la exigencia técnica de las bases (\transv{RT-19.05}{33}), audIT implementará desde el primer mes del contrato un \textbf{Espacio Colaborativo Digital} unificado en la nube de acceso seguro 24/7 para el equipo del proyecto y la contraparte técnica de Curimón. Este repositorio centralizado alojará la documentación formal, especificaciones de diseño, minutas firmadas, registro vivo de riesgos conforme a la norma ISO 31000 (ISO, 2018; \transv{RT-19.04}{33}) y el catálogo de solicitudes de cambio.

El jefe de proyecto tiene dedicación exclusiva durante la implementación y centraliza línea base, cambios, riesgos y seguimiento, conforme a RT-19.01. S7/T-15 dimensiona esa dedicación y los frentes simultáneos; las actas identifican contraparte, responsable de acción y plazo de cierre.

Curimón designará por acta antes del primer sprint a un responsable funcional titular y a un suplente con atribución para ordenar el backlog y aceptar o rechazar incrementos funcionales dentro del alcance aprobado. Se propone responder dudas bloqueantes en dos días hábiles y resolver la revisión de un incremento en cinco días hábiles; la contraparte debe ratificar estos plazos. Si el titular no responde, actúa el suplente; si ninguno decide, el Jefe de Proyecto registra el bloqueo y lo eleva al Comité de Proyecto. La revisión de sprint documenta conformidad de producto, pero la aceptación contractual de H4 y H9 exige las evidencias y actas de S7/S9; el responsable funcional no altera unilateralmente la línea base.

\FloatBarrier
\subsection{Gestión de adquisiciones e integración contractual}\label{sec:6-gestion-adquisiciones}\label{sec:06-gestion-de-adquisiciones-e-integracion-contractual}

La gestión de adquisiciones se estructura para mitigar riesgos de desabastecimiento en la cadena de suministros tecnológicos que pudieran afectar la Ruta Crítica:

\begin{itemize}
\item Adquisición e instalación de componentes vehiculares: La base de T-11 contempla 182 equipos iWave G26I para 148 camiones propios y 34 de terceros sin equipo, más 19 unidades de reposición; contempla 182 lectores CAN CANCrocodile más 19 de reposición. La compra corresponde al CLIENTE. La adquisición, instalación y puesta en servicio se programan con los paquetes y fechas de S7; los 22 equipos previstos para crecimiento de la flota propia se compran durante la operación (T-11; \caso{numeral 14.1}{29}). Esta secuencia sustituye una fecha de compra anticipada no conciliada con el cronograma.
\item Niveles de servicio durante Operación: Los mínimos mensuales de disponibilidad se aplican por criticidad según \art{78.2}{40}: 99,9 \% para servicios críticos y 99,5 \% para servicios de nivel alto; los niveles medio y bajo tienen umbrales propios. El 99,5 \% no es un SLA uniforme para todos los servicios ni reemplaza las metas y mediciones de las etapas de desarrollo y marcha blanca, que se detallan en S9.
\item Conectividad Celular Multicarrier: Contratos de conectividad telemática con SIM card industriales en modalidad APN privada sobre redes de telecomunicaciones de cobertura nacional, con conmutación automática entre operadores para minimizar zonas de silencio.
\end{itemize}

\FloatBarrier
\subsection{Comités de gobernanza, cadencias y mecanismos de decisión}\label{sec:6-gestion-ceremonias}\label{sec:06-comites-de-gobernanza-cadencias-y-mecanismos-de-decision}

El control directivo y operacional del contrato se estructura en estricto apego al marco de gobernanza mandatado por las bases de licitación (\art{71}{37}). La interacción formal entre audIT y Transportes Curimón S.A. se canaliza a través de las cinco instancias que se detallan en la \tabref{06-tabla-6-1}.

\Needspace{20\baselineskip}
\begin{tabla}[ancho=completo, fuente=condensada]{Instancias formales de gobernanza, participantes y cadencias}{06-tabla-6-1}{L{29mm} X X X}
Instancia & Frecuencia & Participantes obligatorios & Propósito y alcance decisional \\ \midrule
Comité Ejecutivo & Mensual & Patrocinador, Gerencia audIT, Administrador del Contrato & Alcance, riesgos mayores \\
Comité de Proyecto & Quincenal & Contraparte Técnica, Jefe de Proyecto & EDT, hitos, bloqueos \\
Comité de Arquitectura & Mensual & Arquitecto, Ciberseguridad, TI Curimón & ADR, seguridad, deuda \\
Comité de Operación & Mensual desde M13 & Operación, Mesa de Ayuda, Contraparte Técnica & SLA, incidentes, marcha blanca \\
Reunión de Seguimiento & Semanal & Equipos de ambas partes & Impedimentos y acciones \\
\end{tabla}

La \tabref{06-tabla-6-1} asigna responsables y cadencias para resolver discrepancias técnicas u operativas. El Comité Ejecutivo decide cambios contractuales y riesgos mayores; el Comité de Proyecto controla la ruta y acuerdos técnicos; el de Arquitectura aprueba ADR y estándares; el de Operación revisa niveles de servicio y problemas recurrentes. Las decisiones, responsables y fechas se conservan en actas; una demora se escala según su impacto. El escalamiento operacional transita desde la Reunión Semanal al Comité de Proyecto ante desviaciones operativas, elevándose al Comité Ejecutivo cuando existe impacto en el alcance, presupuesto o nivel de servicio convenido.

Cuando surge una contingencia o solicitud de cambio que modifique el alcance o los plazos acordados, esta debe tramitarse obligatoriamente mediante el procedimiento de control de cambios regido por el Artículo 72° de las Bases Administrativas (\art{72}{38}). La secuencia de este proceso se detalla en la \tabref{06-tabla-6-2}.

\begin{tabla}[ancho=completo, fuente=condensada]{Procedimiento formal de gestión y control de cambios contractuales}{06-tabla-6-2}{L{29mm} X X X}
Paso & Acción requerida & Responsable formal & Criterio y resultado verificable \\ \midrule
1. Registro & RFC y motivo & Parte solicitante & Solicitud con urgencia \\
2. Análisis & Impacto integral & Jefe de Proyecto y Arquitecto & Informe de impacto \\
3. Revisión & Viabilidad técnica & Comité de Arquitectura & Dictamen y ADR afectado \\
4. Decisión & Aprobar o rechazar & Comité Ejecutivo y ambas partes & Acta y acuerdo escrito \\
5. Ejecución & Rebaselinar y desplegar & Ingeniería y PMO & Sprint o ventana aprobada \\
\end{tabla}

Conforme se establece en la \tabref{06-tabla-6-2}, la RFC registra descripción técnica, justificación operativa y urgencia. El informe de impacto descompone alcance, ruta crítica, costo, riesgo, disponibilidad y pruebas. Arquitectura revisa contratos entre los seis contextos desplegables y seguridad. El acuerdo escrito precede cualquier ejecución, respeta el límite acumulado del 20 \% del valor original y no cambia objeto, garantías mínimas ni hitos obligatorios (\art{72}{38}). La línea base aprobada se actualiza antes del sprint o ventana de mantenimiento. Una alteración sin aprobación previa en acta del Comité Ejecutivo carecerá de validez contractual y no dará derecho a pago adicional (\art{72.5}{38}). Si se suscitaren controversias no resueltas en sede del Comité Ejecutivo en treinta días corridos, operarán los mecanismos de mediación y arbitraje de derecho ante el Centro de Arbitraje y Mediación de Santiago (\art{87}{45}).

\textbf{Compromiso C-01.} audIT formalizará las decisiones y acuerdos de cada sesión de comité en un plazo máximo de veinticuatro horas hábiles en el espacio colaborativo digital.  Métrica: Emisión de minuta formal y registro en espacio digital en menos de 24 horas hábiles tras cada sesión de comité. Se verifica en: Registro de auditoría del repositorio colaborativo. Fuente: compromiso adicional audIT sobre las instancias de \art{71}{37} y el espacio de RT-19.05; las bases no fijan ese plazo de minuta.

\section{Metodología de Desarrollo Software}\label{sec:6-desarrollo}\label{sec:06-metodologia-de-desarrollo-de-software}

El desarrollo de la solución tecnológica se estructura bajo un ciclo de vida evolutivo basado en ingeniería de software guiada por el dominio (Domain-Driven Design, DDD) y prácticas DevSecOps. Se controla deuda técnica, dependencias y tamaño de liberaciones mediante criterios de aceptación verificables.

\FloatBarrier
\subsection{Ciclo de vida adaptado y evolución arquitectónica}\label{sec:6-desarrollo-ciclovida}\label{sec:06-ciclo-de-vida-adaptado-y-evolucion-arquitectonica}

Para conciliar la alta disponibilidad exigida con la continua evolución logística de Transportes Curimón, el desarrollo de software se organiza en torno a los límites de contexto definidos en el Subdocumento 4.1. Este diseño modular desacopla los servicios transaccionales de telemetría de las interfaces de usuario y los algoritmos analíticos.

\begin{itemize}
\item Gestión Evolutiva de Requerimientos: Cada uno de los 42 requerimientos del Formulario T-12 se desglosa en Historias de Usuario documentadas en el repositorio colaborativo. Cada historia incluye criterios de aceptación redactados en formato estructurado (Gherkin: Dado, Cuando, Entonces), sirviendo de especificación viva ejecutable.
\item Diseño de Interfaces API-First: Todo intercambio de datos entre módulos internos y con sistemas legados del cliente se realiza mediante contratos formales: especificación OpenAPI 3.1 para servicios síncronos REST y especificación AsyncAPI 2.6 para eventos telemáticos sobre Event Hubs Premium mediante protocolo Kafka (\art{4.3}{5} y \transv{RT-05.16}{12}).
\item Prevención y Mitigación de Deuda Técnica: Cada sprint de construcción reserva un 15 \% de la capacidad de desarrollo para refactorización, optimización de consultas en base de datos y actualización de librerías base. Se prohíbe la acumulación de advertencias de compilación y se somete el código a inspección estática continua.
\end{itemize}

La \verSeccion{sec:13-innovacion-2} se incorpora como backlog de EDT 12.2, integrado con 4.1/4.2. Sus criterios y métricas se trazan por separado de los 42 requisitos de T-12; no se modifica el significado de un requisito obligatorio para presentarlo como innovación. El piloto y las intervenciones mantienen las ventanas de S7.

\FloatBarrier
\subsection{Ecosistema DevSecOps unificado en GitLab CI Enterprise}\label{sec:6-desarrollo-devsecops}\label{sec:06-ecosistema-devsecops-unificado-en-gitlab-ci-enterprise}

GitLab CI Enterprise centraliza repositorio, incidencias, revisión, ejecución y registro de liberaciones. SonarQube, Trivy, Cosign, herramientas de prueba y Terraform se integran al mismo flujo; usar GitLab no elimina su necesidad ni acredita automáticamente SLSA nivel 3.

\decision{titulo={Orquestación de DevSecOps en GitLab CI Enterprise}, decide={Un solo control de versión y orquestación, con herramientas especializadas integradas}, descarta={Pipelines independientes por equipo con despliegues sin registro común}, criterio={Mantener ID de requisito, commit, prueba, artefacto y despliegue; asumir operación y licencias de las herramientas declaradas}, fuente={\transv{RT-04.03–05}{10}}}

El control de cadena de suministro adopta SLSA v1.2 Build L3: plataforma alojada y endurecida, ejecuciones aisladas entre sí y secretos de firma de procedencia inaccesibles a los pasos definidos por el proyecto. El builder genera la procedencia; un script del propio repositorio que crea una atestación y usa Cosign no basta para L3. Se verifican identidad del builder, parámetros, entradas y digest en la admisión. Es un requisito de implementación a acreditar, no una certificación ya obtenida (SLSA, Build Track \parencite{d3_slsa12}).

El pipeline de entrega continua se ejecuta automáticamente ante cada evento de integración en el repositorio, transitando obligatoriamente por seis fases de validación, ilustradas en la \figref{06-6-2-pipeline}.

\figuraHorizontalAqui{figuras/06-metodologias/6-2-pipeline.pdf}{Fases y puertas de calidad del pipeline DevSecOps unificado}{06-6-2-pipeline}

Como se representa en la \figref{06-6-2-pipeline}, ninguna versión de software puede alcanzar el ambiente de producción sin superar secuencialmente las seis etapas del pipeline. En la \tabref{06-tabla-6-3} se definen los umbrales bloqueantes parametrizados para comprobar los criterios de calidad y seguridad del artefacto.

\begin{tabla}[ancho=completo, fuente=condensada]{Umbrales y políticas de calidad bloqueantes en el pipeline DevSecOps}{06-tabla-6-3}{L{29mm} X X X}
Fase del Pipeline & Herramienta ejecutora & Métrica evaluada & Umbral bloqueante de paso a producción \\ \midrule
1. Pruebas & Runner; dotnet/go/flutter test & Lógica y ramas & $\geq$70 \% / $\geq$80 \%; cero fallos \\
2. Análisis estático & SonarQube & Calidad y seguridad & Gate, A, deuda <5 \%, complejidad $\leq$15 \\
3. Contenedores & Trivy & CVE en imagen & Cero críticos/altos \\
4. Firma y procedencia & Cosign y Syft & Digest y SBOM & Firma y procedencia verificadas \\
5. Infraestructura & Terraform y Checkov & Plan e IaC & Plan revisado; cero bloqueantes \\
6. Análisis dinámico & ZAP y revisión humana & Web y autorización & ASVS L2 trazado; cero críticos/altos \\
\end{tabla}

En la \tabref{06-tabla-6-3}, el 70 \% de lógica de negocio procede de \transv{RT-04.11}{10}; el 80 \% de ramas es meta adicional de audIT y se informa por separado con numerador y denominador. SonarQube exige mantenibilidad A y cero vulnerabilidades críticas/altas, además de los umbrales de deuda y complejidad; Trivy revisa dependencias y capas base. Cada artefacto lleva SBOM CycloneDX, digest, firma y procedencia de un builder aislado, con secretos de firma fuera de los pasos del proyecto. Terraform requiere revisión de plan, escaneo IaC y verificación de identidades, redes y configuración por ambiente. ZAP complementa pruebas manuales de autorización y trazabilidad por control ASVS 4.0 nivel 2; un escaneo no acredita ASVS completo. Los parámetros bloquean automáticamente la promoción cuando no se cumplen. Un fallo aborta el despliegue y notifica al responsable.

\FloatBarrier
\subsection{Infraestructura como código y gestión de configuración}\label{sec:6-desarrollo-iac}\label{sec:06-infraestructura-como-codigo-y-gestion-de-configuracion}

Toda la infraestructura cloud alojada en Microsoft Azure se gestiona bajo el paradigma de Infraestructura como Código (IaC) mediante scripts de Terraform versionados en Git (\transv{RT-03.03}{8}). Se implementa un modelo de entornos rigurosamente aislados:

\begin{itemize}
\item Segregación de Ambientes: DEV, QA, PREPROD, PROD y DR separados, con suscripciones/grupos de recursos, redes, identidades y almacenes según la topología de S4.2; diferencias de PREPROD frente a PROD declaradas, con controles y pruebas que impiden acceso no autorizado. Datos sintéticos en DEV/QA; una seudonimización reversible no se denomina anonimización.
\item Aprovisionamiento Inmutable: Los servidores y clusters de Azure Kubernetes Service (AKS) no admiten modificaciones manuales en caliente vía consola. Las políticas y límites de escalamiento se registran como código; el autoscaler ejecuta esas políticas sin requerir un commit por cada réplica. Cualquier cambio manual excepcional documentado o modificación de política debe quedar auditado, someterse a revisión por pares y desplegarse mediante pipeline.
\item Gestión Centralizada de Secretos: Ninguna credencial, clave privada o cadena de conexión se almacena en el código fuente. Se utiliza Azure Key Vault integrado con identidades administradas (Managed Identities) y rotación programada automática.
\end{itemize}

\FloatBarrier
\subsection{Ceremonias técnicas, artefactos y criterios de salida}\label{sec:6-desarrollo-ceremonias}\label{sec:06-ceremonias-tecnicas-artefactos-y-criterios-de-salida}

El trabajo colaborativo del equipo de desarrollo se organiza en torno a un flujo de trabajo de control de versiones GitFlow adaptado. Las ramas principales son:

\begin{itemize}
\item Rama \texttt{main}: Contiene versiones estables liberables y está protegida. Se exige firma criptográfica de commits y dos aprobaciones independientes del autor antes de la liberación; producción se despliega automáticamente. El registro de despliegue identifica commit, etiqueta, digest, configuración y entorno efectivo; fusionar a main no equivale a haber desplegado.
\item Rama \texttt{develop}: Rama de integración continua donde convergen las nuevas funcionalidades validadas mediante pruebas unitarias.
\item \textbf{Ramas temáticas:} Ramas de trabajo de corta duración, sujetas obligatoriamente a revisión por pares mediante Merge Requests. Las normas ISO citadas en otros capítulos sustentan arquitectura, calidad de datos y riesgos; no se presentan como fuente del flujo de ramas.
\end{itemize}

La \tabref{06-tabla-6-4} define las condiciones para iniciar una historia y aceptar un incremento de software mediante Definition of Ready (DoR) y Definition of Done (DoD).

\begin{tabla}[ancho=completo, fuente=condensada]{Criterios formales de salida: Definición de Preparado y Definición de Terminado}{06-tabla-6-4}{L{36mm} X X}
Nivel de control & Artefacto / Fase evaluada & Criterios de aceptación obligatorios \\ \midrule
Definición de Preparado (DoR) & Historia de Usuario / Requerimiento funcional & Requerimiento trazado unívocamente al Formulario T-12. Criterios de aceptación definidos en Gherkin. Dependencias técnicas resueltas. Mockups de interfaz aprobados. \\
Revisión por Pares & Merge Request (MR) en GitLab & Al menos un revisor independiente para integración; dos aprobaciones independientes del autor antes de liberar a main/producción. Verificación de adherencia a guías de estilo, comentarios arquitectónicos y ausencia de duplicidad de código. \\
Definición de Terminado (DoD) & Incremento de Software / Release Candidate & Código fusionado en rama objetivo. Cobertura automatizada de lógica de negocio $\geq$ 70 \% (\transv{RT-04.11}{10}) y cobertura de ramas $\geq$ 80 \% como meta de audIT, con numeradores, denominadores y exclusiones reportados por separado; no se promedian ni sustituyen entre sí. Quality Gate SonarQube superado. Contenedor firmado con Cosign y registrado en Azure Container Registry con SBOM. Documentación OpenAPI actualizada. Despliegue exitoso en staging sin regresiones operativas. \\
\end{tabla}

La aplicación de los criterios de la \tabref{06-tabla-6-4} asegura que cada módulo tecnológico entregado a Transportes Curimón cumpla con los estándares industriales de mantenibilidad, solidez arquitectónica y ciberseguridad exigidos por las bases.

\textbf{Compromiso C-02.} audIT mantendrá un tiempo de despliegue continuo en ambiente de pruebas inferior a quince minutos para incrementos de software validados.  Métrica: Despliegue automatizado en staging en menos de 15 minutos tras aprobación de Merge Request. Se verifica en: Métricas de ejecución del pipeline en GitLab CI Enterprise. Fuente: objetivo adicional audIT, trazado a FEP02 RT-04.12 y RT-03.03; no es plazo impuesto por las bases.

\FloatBarrier
\subsection{Cadencia, entrega y reversión verificable}\label{sec:06-cadencia-entrega-y-reversion-verificable}

Cada sprint de dos semanas incluye planificación con objetivo y criterios, reunión diaria de quince minutos, refinamiento de dependencias, revisión con representantes de torre/conductores/transportistas y retrospectiva. Se conservan backlog, diseño de interfaz, ADR, contratos, código, pruebas, SBOM y decisión de liberación. El Comité de Proyecto quincenal recibe avance por paquete, riesgos y desviaciones; la revisión de sprint es validación de producto y no reemplaza aceptación contractual.

El plan inicial de los ocho sprints A09 y cinco A20 se adjunta en T-10, con requisitos y prueba por incremento. Para cada sprint, la capacidad neta de construcción se calcula como días-persona disponibles multiplicados por el factor de foco de S7 y por 0,85 tras reservar 15 \% para deuda técnica; vacaciones, guardias, inducción y otras ausencias se restan antes de aplicar ese factor. S7 fija 70 \% para E2: su capacidad neta inicial equivale a 59,5 \% de los días-persona disponibles. No se traduce ese porcentaje en horas contratadas sin el calendario y dedicaciones de T-15. El responsable funcional titular o suplente revisa el resultado con el perfil usuario indicado en T-10; QA registra pruebas y defectos. Una historia sin criterio de aceptación o dependencia habilitante no entra al sprint.

Durante M21–M56, Kanban reserva un carril visible de deuda técnica con entrada por defecto, riesgo, costo de postergación y prueba de regresión. El Líder de Operación acuerda mensualmente con el Jefe de Proyecto la capacidad en horas y personas a partir de la dotación efectiva de S11/T-15; el 15 \% de Scrum no se traslada automáticamente a operación. Hasta que esa capacidad esté dimensionada y valorizada, el carril es una decisión metodológica y no una asignación de esfuerzo ya acreditada.

DEV, QA, PREPROD, PROD y DR deben estar operativos en H3 (\transv{RT-04.01}{10}). La rama main identifica lo liberado; develop integra incrementos y las ramas temáticas son breves. Se construye el artefacto una sola vez; su digest, firma y configuración externa permiten promoción QA $\rightarrow$ PREPROD $\rightarrow$ PROD sin recompilar (RT-04.08). DAST requiere despliegue temporal en QA/PREPROD, aunque sea la sexta puerta antes de producción. Una comprobación que exige HIL, prueba de terreno o revisión humana se registra como evidencia adicional, sin simularla dentro del pipeline.

Se elige canario para servicios y actualización por unidad para equipos a bordo. Antes de cada pase se comprueban métricas, reversa, esquema compatible y preservación de evidencia. La reversión automática retorna a la versión previa compatible, mantiene eventos no confirmados y conserva documentos emitidos; después reconcilia diferencias. Las migraciones de esquema siguen expandir–migrar–contraer y son versionadas/reversibles (RT-04.10). Contraer columnas o transformar datos de forma irreversible se posterga hasta que la versión anterior deje de necesitarlas; no se promete que toda transformación se deshaga con un rollback de contenedor.

La \tabref{06-tabla-6-5} declara metas de entrega para RT-04.12 y conserva el máximo contractual de cambios fallidos del Artículo 78.3; se miden por contexto en operación y se ajusta capacidad si fallan, sin alterar mínimos de seguridad ni congelamientos de S7.

\begin{tabla}[ancho=completo, fuente=condensada]{Indicadores de entrega y restauración}{06-tabla-6-5}{L{36mm} X X}
Indicador & Meta propuesta & Evidencia \\ \midrule
Frecuencia de despliegue & Incremento disponible cada sprint; producción sólo en ventanas aprobadas sin congelamiento & Historial de liberaciones por servicio y calendario S7 \\
Tiempo desde commit aceptado hasta producción & $\leq$5 días hábiles fuera de congelamientos, con pruebas/gates completos & Marcas de commit, aceptación y despliegue; exclusiones justificadas \\
Cambios fallidos en producción & $\leq$5 \% de los despliegues del mes (FEP01, Art. 78.3, p. 40) & Despliegues con fallo, reversa o incidente atribuible / total de despliegues en producción del mes, ligados al digest \\
Restauración tras cambio fallido & Objetivo $\leq$30 min para reversa de aplicación, sin sustituir SLA de incidente ni RTO de DR & Cronómetro de detección, reversa y verificación de integridad \\
\end{tabla}

La \tabref{06-tabla-6-5} distingue despliegue en pruebas (<15 min propuesto en C-02), lead time, reversa de una aplicación y recuperación completa de un sitio. Ninguno se infiere del éxito del pipeline. Los objetivos se verifican con S9; la aceptación E1/E2 mantiene sus marchas blancas y el régimen crítico 99,9 \% del Artículo 78.2.

\FloatBarrier
\subsection{Puerta de salida G13 y funciones de sprint 7}\label{sec:06-liberacion-g13}

G13 libera la versión de entrenamiento probada y congelada en PREPROD, con datos ficticios y material por perfil. Para habilitar los diez días iniciales de familiarización de A13, el candidato debe incluir los cinco flujos de la \tabref{06-tabla-6-6-g13}. Esta definición complementa la línea base de capacitación del Subdocumento 7 y conserva la duración de A13 y su adelanto condicionado. Los casos trazados en el Formulario T-17 se ejecutan sobre el candidato de PREPROD antes de liberar G13.

\Needspace{20\baselineskip}
\begin{tabla}[ancho=completo, fuente=condensada]{Funciones mínimas de familiarización y criterios previstos para liberar G13}{06-tabla-6-6-g13}{L{31mm} L{24mm} L{31mm} X}
Función & Requisito T-12 & Casos T-17 & Criterio previsto para G13 \\ \midrule
Acceso por rol y privacidad & RF-022, RNF-013 & CP-SEC-04/05 & Acceso/denegación auditados \\
Asignación elegible & RF-001, RF-002 & CP-UNIT-01/02; CP-SYS-01/05 & Reserva válida; rechazo inválido \\
Jornada y descanso & RF-003 & CP-INT-03; CP-SYS-05 & Decisión y motivo registrados \\
Alerta de conducción & RF-027 & CP-UNIT-18 & Aviso antes de cinco horas \\
Conectividad intermitente & RNF-002 & CP-UNIT-23 & Evento íntegro, sin duplicado \\
\end{tabla}

La \tabref{06-tabla-6-6-g13} limita el certificado de familiarización a cinco flujos de software con datos ficticios. En acceso, el rol permitido entra, el ajeno es rechazado y concesión/revocación quedan auditadas. En asignación, una vigencia inválida o una segunda reserva activa bloquea. La prueba de descanso usa escenarios sintéticos de ocho horas y seis horas en 24 horas; registra decisión y causa. La alerta reproduce 4 h 15 min de conducción y un lugar seguro alcanzable en 35 min; debe avisar antes de cinco horas con diez minutos de margen. El corte de enlace incluye reinicio y reenvío doble, conservando un evento íntegro sin duplicarlo. G13 se cumple cuando todos pasan en el mismo candidato de PREPROD destinado a capacitación y no quedan defectos bloqueantes abiertos.

El equipo de desarrollo conserva el identificador y digest del artefacto, configuración de PREPROD, versión del conjunto ficticio, informes de ejecución y reejecución de defectos, materiales por perfil y acta de congelamiento. El equipo de aseguramiento de calidad revisa la trazabilidad de la evidencia. Un cambio posterior al congelamiento obliga a repetir los casos afectados. CP-HW-03/06/14, la autonomía de 72 h y la validación en ruta requieren evidencia propia en S9/T-17; la simulación en PREPROD no las sustituye.

G13 habilita los diez días iniciales de familiarización; los flujos completos, las evaluaciones y el complemento del sprint 8 se mantienen en los treinta días restantes de A13.

\FloatBarrier
\subsection{Evidencia asociada a T-9 y T-10}\label{sec:06-evidencia-asociada-a-t-9-y-t-10}

El \form{T-9}{61} remite a 6.1, interesados, comités, adquisiciones, cambios y calendario; el \form{T-10}{61} remite a 6.2, ciclo de vida, herramientas, ambientes, criterios, cadencias y reversión. Los anexos son índices de evidencia del capítulo, no inventarios que reemplacen el análisis.
